\section{Eliminating a shared reflection prefix}
\label{sec:v77-prefix}

We first record the finite variational facts, including the sign and
nondegeneracy needed to compare two actual experiments.  No periodicity
is used in this section.

\begin{lemma}[Regular finite actions]\label{lem:v77-regular-action}
Let a finite trajectory in a dispersing table have clear flights of
positive length and nongrazing reflections on positive-curvature
boundaries.  Use arbitrary regular scalar coordinates on its successive
boundary patches.  After restricting the patches, its stationary
trajectory is uniquely determined near the reference trajectory by the
two endpoint coordinates.  Its length $W$ is smooth and $W_{st}\ne0$.
The corresponding initial and final covectors satisfy
\[
 p=-W_s(s,t),\qquad q=W_t(s,t).
\]
All these constructions and their fixed-order derivatives depend
continuously on the boundary embeddings on smaller closed domains.
\end{lemma}
\begin{proof}
Write $r_i=R_i(x_i)$, $e_i=(r_{i+1}-r_i)/\ell_i$ and
$\ell_i=|r_{i+1}-r_i|$.  For a variation $v_i=R_i'(x_i)h_i$, the
second differential of the $i$th chord is
\[
 \ell_i^{-1}|P_{e_i^\perp}(v_{i+1}-v_i)|^2
 +e_i\cdot\{R_{i+1}''h_{i+1}^2-R_i''h_i^2\}.
\]
At an interior reflection the coefficient from the second terms is
$(e_{i-1}-e_i)\cdot R_i''h_i^2$.  With outward obstacle normal $n_i$,
reflection gives $e_{i-1}-e_i=-2\nu_i n_i$, $\nu_i>0$, and
$n_i\cdot R_i''=-\kappa_i|R_i'|^2$.  Thus, with endpoints fixed,
\begin{equation}\label{eq:v77-hessian}
 \delta^2\sum_i\ell_i
 =\sum_i\ell_i^{-1}|P_{e_i^\perp}(v_{i+1}-v_i)|^2
  +\sum_{i\text{ interior}}2\kappa_i\nu_i|R_i'|^2h_i^2>0
\end{equation}
for every nonzero interior variation.  Tangential terms caused by the
choice of parameter vanish by stationarity.  The interior Hessian is
therefore positive definite.  The implicit function theorem constructs
the unique nearby stationary path; continuity preserves clearance and
reflection, hence it remains a physical path.

The mixed derivative of an individual chord is
\begin{equation}\label{eq:v77-edge-twist}
 (\ell_i)_{x_ix_{i+1}}
 =-\ell_i^{-1}
   (e_i^\perp\cdot R_i')(e_i^\perp\cdot R_{i+1}').
\end{equation}
Both factors are nonzero by nongrazing.  The interior Hessian is
tridiagonal.  If there are interior sites, the mixed derivative of the
stationary value is the negative product of its two endpoint couplings
and the opposite-corner entry of the inverse interior Hessian.  The
cofactor of that corner is, up to sign, the product of all intervening
off-diagonal entries.  Each factor is nonzero by
\eqref{eq:v77-edge-twist}, and the determinant is positive by
\eqref{eq:v77-hessian}.  Hence $W_{st}\ne0$.  With no interior site,
this is just \eqref{eq:v77-edge-twist}.

The first variation contains only the endpoint terms, giving the two
covector formulas.  The implicit-function and cofactor arguments apply
with parameters and can be differentiated at any fixed order.  On a
compact strictly regular domain their inverses have finite bounds.
\end{proof}

Suppose $\gamma$ goes from $I$ to $J$ and its extension $\gamma e$
goes from $I$ to $K$.  Write
\[
 U(s,t)=W_\gamma(s,t),\quad V(s,v)=W_{\gamma e}(s,v),\quad
 \ell(t,v)=|R_K(v)-R_J(t)|.
\]
The stationary composition identity on the common physical branch is
\begin{equation}\label{eq:v77-stationary-composition}
 V(s,v)=U(s,t(s,v))+\ell(t(s,v),v),\qquad
 U_t(s,t(s,v))+\ell_t(t(s,v),v)=0.
\end{equation}
The covector is unchanged by reflection at $J$.  This explains why the
same stationarity equation joins the prefix to its last chord even
though the prefix record stops at $J$.

\begin{proposition}[Nonlinear prefix cancellation]
\label{prop:v77-cancellation}
Near any such extended regular trajectory there are closed coordinate
intervals $I_0,J_0,K_0$ on which the equation
\begin{equation}\label{eq:v77-matching}
 H(s,t,v):=V_s(s,v)-U_s(s,t)=0
\end{equation}
has a unique solution $s=\sigma(t,v)\in\operatorname{int}I_0$ for
$(t,v)\in J_0\times K_0$.  The domains can be chosen with
$|H_s|\ge\mu>0$ and a strict endpoint sign bracket.  The absolute
last-flight distance is then determined by the two actions:
\begin{equation}\label{eq:v77-distance-cancellation}
 \ell(t,v)=V(\sigma(t,v),v)-U(\sigma(t,v),t).
\end{equation}
There is no additive ambiguity.  On bounded classes with these margins,
for every $k\ge0$ the map $(U,V)\mapsto\ell$ is locally Lipschitz
from $C^{k+2}$ on slightly larger fixed domains to $C^k(J_0\times K_0)$.
\end{proposition}
\begin{proof}
For a fixed initial state, the common prefix in the two experiments is
the same actual path.  Eliminating the interior coordinates of the
prefix in the full second variation leaves a strictly positive Schur
complement
\[
 B=U_{tt}+\ell_{tt}>0
\]
at the matched path, by \eqref{eq:v77-hessian}.  Differentiate the
stationarity equation in \eqref{eq:v77-stationary-composition} at fixed
$v$.  It gives $t_s=-U_{st}/B$.  Also $V_s=U_s(s,t(s,v))$, so at a
matched triple
\begin{equation}\label{eq:v77-matching-derivative}
 H_s=V_{ss}-U_{ss}=U_{st}t_s=-\frac{U_{st}^2}{B}<0.
\end{equation}
The implicit function theorem therefore solves for $s$ as a function
of $(t,v)$.  Restrict to a product box about the reference triple on
which $H_s<0$.  Choose its two $s$ endpoints on opposite sides of the
root and then shrink $J_0,K_0$ so that the endpoint signs are uniform.
Compactness gives the stated strict margins.  Monotonicity gives
uniqueness on this entire specified box, not just formal uniqueness of
a derivative at the reference point.

Equation \eqref{eq:v77-matching} equates initial covectors at the same
initial scalar label.  A nongrazing ray is determined by this covector
and the choice of the outgoing half-plane.  The branch labels select
that half-plane and the same first reflection.  Determinism and local
uniqueness of the physical path identify the full common prefix.
Subtracting its length proves \eqref{eq:v77-distance-cancellation}.
Equivalently the same conclusion follows by substituting the root in
\eqref{eq:v77-stationary-composition}.

For stability, retain the strict root bracket and $|H_s|\ge\mu$ on a
slightly larger compact box.  For a nearby pair $(\widetilde U,
\widetilde V)$ these margins persist.  The mean value theorem first
gives
\[
 \|\widetilde\sigma-\sigma\|_{C^0}
 \le C\|\widetilde H-H\|_{C^0}.
\]
Differentiate $H(\sigma(t,v),t,v)=0$.  At every order, the highest
derivative of $\sigma$ is multiplied by $H_s$; the remaining terms
involve lower derivatives and derivatives of $H$.  Subtracting the two
equations and inducting gives
\[
 \|\widetilde\sigma-\sigma\|_{C^k}
 \le C\|\widetilde H-H\|_{C^{k+1}}
 \le C\bigl(\|\widetilde U-U\|_{C^{k+2}}
             +\|\widetilde V-V\|_{C^{k+2}}\bigr).
\]
The extra derivative controls composition at shifted roots.  Products,
composition on compact boxes and the mean value theorem applied to
\eqref{eq:v77-distance-cancellation} prove the claimed bound.  The
order two loss is a sufficient bound, not an optimality assertion.
\end{proof}

In symplectic notation this argument is
$\mathcal F_e=\mathcal F_{\gamma e}\circ\mathcal F_\gamma^{-1}$,
where $\mathcal F_W:(s,-W_s)\mapsto(t,W_t)$.  The explicit root
formula also recovers the additive constant in its generating function.
It does not use any measured momentum: derivatives of recovered actions
supply all covectors in the argument.  Different source weights for the
two branches were already eliminated separately by their clock triples.

The local root bracket is a witness used in constructing a scalar
gate.  It need not be supplied to the observer: the next section
recovers the matching domain from the actions.  An arbitrary choice
among roots of unrelated stationary functions is not justified.
Formula \eqref{eq:v77-matching-derivative} proves that suitable brackets
exist for every genuine shared-prefix pair.  The finite-atlas theorem
uses precisely these physical boxes.
